% 19.11.2025; 17:11:50


\section{Zusammenfassung}%
\begin{frame}%
    \frametitle{Zusammenfassung}
    {\large Ein Kompromiss:}

    \begin{outline}
        \1 Objektorientierte Programmierung
            \2 State sollte in Klassen gekapselt sein
            \2 Globalen State vermeiden
        \vspace{.3em}

        \1 Deklarative Programmierung
            \2 State allgemein vermeiden
            \2 State-Änderungen vermeiden
            \2 Deklarative-Primitives nutzen
            \2 Immutability > Datenkapselung
        \vspace{.3em}

        \1 Design Pattern
            \2 Bieten Struktur, Standardisierung und Namenskonvention
    \end{outline}

\end{frame}




\note {
    \begin{outline}
        \1 Am Ende kommt es auf die State Frage an
        \1 Objektorientierte Programmierung
            \2 State sollte in Klassen gekapselt sein
            \2 Globaler State ist zu vermeiden
        \1 Deklarative Programmierung
            \2 State-Änderungen / Änderbarer State ist zu vermeiden
            \2 Deklarative-Primitives nutzen
                \3 Funktionen wie reduce, forEach, etc. werden von Sprache idR. bereits bereitgestellt
            \2 Zu starke Gewichtung dieser Ideale in den meisten Kontexten schwierig, etwas Deklarative Programmierung macht Code häufig besser
            \2 Immutability > Datenkapselung
        \1 Design Pattern
            \2 Bieten Struktur, Standardisierung und Namenskonvention
    \end{outline}



}

